%%\documentclass[sn-nature]{sn-jnl}% Style for submissions to Nature Portfolio journals
%%\documentclass[sn-basic]{sn-jnl}% Basic Springer Nature Reference Style/Chemistry Reference Style
\documentclass[sn-mathphys-num]{sn-jnl}% Math and Physical Sciences Numbered Reference Style 
%%\documentclass[sn-mathphys-ay]{sn-jnl}% Math and Physical Sciences Author Year Reference Style
%%\documentclass[sn-aps]{sn-jnl}% American Physical Society (APS) Reference Style
%%\documentclass[sn-vancouver,Numbered]{sn-jnl}% Vancouver Reference Style
%%\documentclass[sn-apa]{sn-jnl}% APA Reference Style 
%%\documentclass[sn-chicago]{sn-jnl}% Chicago-based Humanities Reference Style

\usepackage{graphicx}%
\usepackage{multirow}%
\usepackage{amsmath,amssymb,amsfonts}%
\usepackage{amsthm}%
\usepackage{mathrsfs}%
\usepackage[title]{appendix}%
\usepackage{xcolor}%
\usepackage{textcomp}%
\usepackage{manyfoot}%
\usepackage{booktabs}%
\usepackage{algorithm}%
\usepackage{algorithmicx}%
\usepackage{algpseudocode}%
\usepackage{listings}%

\theoremstyle{thmstyleone}%
\newtheorem{theorem}{Theorem}%  meant for continuous numbers
%%\newtheorem{theorem}{Theorem}[section]% meant for sectionwise numbers
%% optional argument [theorem] produces theorem numbering sequence instead of independent numbers for Proposition
\newtheorem{proposition}[theorem]{Proposition}% 
%%\newtheorem{proposition}{Proposition}% to get separate numbers for theorem and proposition etc.

\theoremstyle{thmstyletwo}%
\newtheorem{example}{Example}%
\newtheorem{remark}{Remark}%

\theoremstyle{thmstylethree}%
\newtheorem{definition}{Definition}%

\raggedbottom
%%\unnumbered% uncomment this for unnumbered level heads

\begin{document}

\title{Final paper}

%%=============================================================%%
%% GivenName	-> \fnm{Joergen W.}
%% Particle	-> \spfx{van der} -> surname prefix
%% FamilyName	-> \sur{Ploeg}
%% Suffix	-> \sfx{IV}
%% \author*[1,2]{\fnm{Joergen W.} \spfx{van der} \sur{Ploeg} 
%%  \sfx{IV}}\email{iauthor@gmail.com}
%%=============================================================%%

\author[1]{\fnm{Annie} \sur{Dai}}\email{anniedai@umd.com}
\equalcont{These authors contributed equally to this work.}

\author[1]{\fnm{Emma} \sur{Shroyer}}\email{eshroyer@gmail.com}
\equalcont{These authors contributed equally to this work.}

\author[1]{\fnm{Nathalie} \sur{Bonin}}\email{nbonin@gmail.com}
\equalcont{These authors contributed equally to this work.}

\affil[1]{\orgdiv{Computer Science Department}, \orgname{University of Maryland}, \orgaddress{\city{College Park}, \state{Maryland}, \country{USA}}}

%%==================================%%
%% Sample for unstructured abstract %%
%%==================================%%

\abstract{The abstract serves both as a general introduction to the topic and as a brief, non-technical summary of the main results and their implications. Authors are advised to check the author instructions for the journal they are submitting to for word limits and if structural elements like subheadings, citations, or equations are permitted.}
Salmon provides a fast and accurate way to estimate transcript abundance. This method uses a dual-phase inference algorithm and takes into account fragment biases for transcript verification. However, there are some areas where Salmon could improve. The inference algorithms it uses, EM and VBEM, are only point estimates and cannot be checked if there is no ground truth available. Further, the choice of prior when using the VBEM algorithm can affect the accuracy of the results. In this work, we test how averaging the inference model predictions affect the accuracy of the inference estimation during Salmon’s offline phase. 
%%================================%%
%% Sample for structured abstract %%
%%================================%%

% \abstract{\textbf{Purpose:} The abstract serves both as a general introduction to the topic and as a brief, non-technical summary of the main results and their implications. The abstract must not include subheadings (unless expressly permitted in the journal's Instructions to Authors), equations or citations. As a guide the abstract should not exceed 200 words. Most journals do not set a hard limit however authors are advised to check the author instructions for the journal they are submitting to.
% 
% \textbf{Methods:} The abstract serves both as a general introduction to the topic and as a brief, non-technical summary of the main results and their implications. The abstract must not include subheadings (unless expressly permitted in the journal's Instructions to Authors), equations or citations. As a guide the abstract should not exceed 200 words. Most journals do not set a hard limit however authors are advised to check the author instructions for the journal they are submitting to.
% 
% \textbf{Results:} The abstract serves both as a general introduction to the topic and as a brief, non-technical summary of the main results and their implications. The abstract must not include subheadings (unless expressly permitted in the journal's Instructions to Authors), equations or citations. As a guide the abstract should not exceed 200 words. Most journals do not set a hard limit however authors are advised to check the author instructions for the journal they are submitting to.
% 
% \textbf{Conclusion:} The abstract serves both as a general introduction to the topic and as a brief, non-technical summary of the main results and their implications. The abstract must not include subheadings (unless expressly permitted in the journal's Instructions to Authors), equations or citations. As a guide the abstract should not exceed 200 words. Most journals do not set a hard limit however authors are advised to check the author instructions for the journal they are submitting to.}

\keywords{Salmon, Tuna, Onefish, Twofish, Redfish, Bluefish}

%%\pacs[JEL Classification]{D8, H51}

%%\pacs[MSC Classification]{35A01, 65L10, 65L12, 65L20, 65L70}

\maketitle
\section{Background}
A key characteristic of many newly-formed eukaryotic RNA sequences is the 
inclusion of non-coding and coding regions, which are often referred to as introns and exons. 
Prior to translation, transcript maturation must occur via splicing, 
a process that revolves around the displacement of introns from the mRNA sequence. 
However, depending on environmental conditions and sequence variations, 
some exons are often removed from the mRNA as well, eventually resulting in 
the creation of a different protein isoform. 
Unfortunately, the many factors related to alternative splicing are still poorly understood 
for most eukaryotic proteins, hence the reasoning behind the continuation of transcriptome projects.
\\
An essential facet to transcriptomics is the sequencing techniques used to generate transcript 
fragments for data analysis. Prior approaches to transcriptome quantification involved the use 
of microarrays, which has important limitations such as relatively high risks of cross-hybridization, 
low range of detection, and complicated normalization methods for differential expression analysis 
\cite{Wang_Gerstein_Snyder_2009}. 
RNA-Seq, on the other hand, doesn’t face such limitations and retains microarrays’ 
advantage over classical sequencing such as high-throughput capabilities and the use of inexpensive technology. 
It has been revolutionary in furthering the field of transcriptomics, specifically with regards to eukaryotic cell transcriptomes. 
In particular, recent RNA-seq experiments yielded meaningful results in quantifying expression 
difference across different tissues \cite{Glinos_et_al._2022}, distinguishing host-pathogen interactions 
\cite{Pisu_Huang_Grenier_Russell_2020}, and identifying the effect of structural variants on splicing regulators 
\cite{Pascal_et_al._2023}. 
\\	
Equally important to transcriptomics is abundance quantification. 
Estimations can be done through an alignment to a reference transcriptome sequence; nevertheless, a constant struggle with quantifying eukaryotic samples is to map read fragments back to the original transcript as opposed to the original. Although the detection of protein family abundances can be useful on its own, oftentimes researchers are more so interested in transcript abundances. Since isoforms of the same protein will retain common exons, identifying the source of each read is a tedious process that requires complex statistical methods.
In the past, Salmon was developed to address this particular issue. 
Given a known transcript and a set of sequenced fragments, it quantifies the relative abundance of each transcript in the sample using a dual-phase inference procedure. Salmon improves accuracy compared to other models, because it takes into account sample-specific biases to improve accuracy. When these biases are not accounted for, calculations like the false discovery rate cannot be controlled for. Using multiple inference steps, Salmon improves its abundance estimates using either the VBEM or EM inference algorithms. However, both methods have drawbacks. Both algorithms return point estimates of abundances, and we are uncertain about the returned estimates without ground truth. Further, the choice of  Bayesian priors used in the VBEM algorithm also has considerations. A small prior leads to sparser results than EM. However, a larger prior may result in more estimated non-zero abundance than EM. Prior simulated tests [Salmon documentation] show that VBEM with a small prior can lead to more accurate estimates. However, there has been much research into improving the results of these algorithms such as model averaging. \cite{doi:10.1177/2515245919898657}\cite{11c48783-68bb-36b7-8053-175211b0eaa8}\cite{hoeting_bayesian_1999} Given the trade offs of the two algorithms, we sought to know if the averaged prediction for EM and VBEM inference is better than one algorithm or another. Further, if this ensemble estimate is better, could it be improved by varying the Bayesian priors.
We have found that the average of these outputs [tie in results]

\subsection{Salmon}
Salmon uses a dual-phase inference procedure to provide fast and accurate estimates. 
In an online and offline step, Salmon is able to calculate and improve the abundance 
estimates for a transcript based on fragment GC-content and positional biases. 
Salmon initially utilizes raw reads in a quasi-mapping phase where it performs direct 
quantification instead of enacting a traditional alignment. 
Salmon then moves to an online inference phase. 
Here, Salmon uses a variant of stochastic, collapsed variational Bayesian inference to solve 
a variation Bayesian inference problem. 
In this phase, Salmon estimates the initial expression levels, auxiliary parameters, 
the ‘foreground’ bias models, and fragment equivalence classes.In the offline phase, 
Salmon improves the estimates it calculated in the online phase. Here, the user can 
specify the inference algorithm of EM or VBEM. The chosen algorithm runs to convergence 
on this data, outputting an optimized array of abundances. 
Optionally, using the converged abundances and fragment equivalence classes, 
Salmon can also draw and save estimates from the posterior distribution using Gibbs or bootstrap sampling. \cite{patro_salmon_2017}
\subsection{EM \& VBEM}

The Expectation Maximization (EM) algorithm optimizes the likelihood of the parameters given the data. This algorithm returns a point estimate of the abundances. 
This algorithm is  the default of Salmon.
\\
The Variational Bayesian Expectation Maximization (VBEM) algorithm accounts for the sparsity of the data. It takes a Bayesian nucleotide prior that controls for the sparsity of the data. The default prior for Salmon was $1 \times 10 ^{-2}$.

Each transcriptome \(\mathcal{T}\) is composed of transcripts \(t\).
Each transcript is a nucleotide sequence which can be described through
its length \(l\), effective length \(\tilde l\), and its count
\(c\)---which is the number of times that \(t\) occurs in a given
sample. There are \(M\) transcripts in a given transcriptome.

The probability that a given fragment originates from a transcript \(t\)
depends on the length of that transcript relative to all other
transcripts in the transcriptome. We define this nucleotide fraction
\(\eta\) for a given transcript \(t\) as

\[\eta_i= \dfrac{c_i \cdot \tilde l _i}{\sum^{M}_{\pmb{j} = 1}c_j \cdot \tilde l _j}\]

We obtain a transcript fraction \(\tau\) for a given transcript \(t\) by
normalizing its nucleotide fraction against the effective length of all
transcripts.

\[\eta_i= \dfrac{c_i \cdot \tilde l _i}{\sum^{M}_{\pmb{j} = 1}c_j \cdot \tilde l _j}\]

Let's say that the true nucleotide fraction for a transcript \(t\) is
\(\pmb{\eta}\). We can describe the probability of observing a set of
sequenced fragments \(\mathcal{F}\) as a

\[
$$	
\begin{align*}
	\Pr\left( \mathcal{F}| \pmb{\eta}, Z, \mathcal{T} \right) &= \prod^{N}_{j=1} \Pr\left( f_j| \pmb{\eta}, Z, \mathcal{T} \right)\\
	&= \prod^{N}_{j=1}\sum^{M}_{i=1} \Pr\left( t_i| \pmb{\eta} \cdot \Pr\left( f_j|t_i,z _{ij} =1 \right)  \right)
\end{align*}
$$
\]

where \(N\) is the number of fragments in \(\mathcal{F}\), \(Z\) is a
relationship matrix where \(z _{ij} =1\) when fragment \(f_j\) is
derived from \(t_i\).

We want to obtain \(\pmb{\alpha}\), which is the estimated number of
reads from each transcript. We describe the maximum likelihood estimates
as:

\[\mathcal{L} \left( \pmb{\alpha} |\mathcal{F}, \pmb{Z},\mathcal{T} \right)  = \prod^{N}_{j=1} \sum^{M}_{i=1} \hat \eta_i \Pr\left(  f_j|t_i\right)\]

Written in terms of equivalence classes \(\pmb{C}\)

\[\mathcal{L} \left( \pmb{\alpha} |\mathcal{F}, \pmb{Z},\mathcal{T} \right)  = \prod^{}_{\mathcal{C}^j \in \pmb{C}} \left( \sum_{t_i \in  \pmb{t}^j} \hat \eta_i w_t^j\right) ^{d^{j}}\]

The abundances \(\pmb{\hat \eta}\) are computed directly from
\(\pmb{\alpha}\)

\[\hat \eta_i= \dfrac{\alpha_i}{\sum_j \alpha_j}\]

Then we apply an update function

\[\alpha_i ^{u+1} = \sum^{}_{\mathcal{C}^j \in \pmb{\mathcal{C}}} d ^{j} \left( \dfrac{\alpha_i^u w_i^j}{\sum^{}_{t_k \in \pmb{t}}} j \alpha^u_k w_k^j\right)\]

Until the maximum relative difference in \( \pmb{\alpha}\) is

\[\Delta \left( \alpha^u, \alpha ^{u+1} \right) = \mathrm{max} \dfrac{\left| a_i^u  - \alpha_i ^{u+1}\right| }{ \alpha_i ^{u+1}} < 1 \times 10 ^{-2}\]

for all \(\alpha_i ^{u+1} > 1 \times 10 ^{-8}\), at which point we
derive the estimated nucleotide fraction

\[\hat \eta_i =\dfrac{\alpha_i ^{\prime} }{ \sum_j \alpha ^{\prime} _j}\]

Variational Bayes Optimization

Optionally, the we can apply variational bayeseian optimization where
the update function is

\[\alpha_i ^{u+1} = \sum^{}_{\mathcal{C}^j \in \pmb{\mathcal{C}}} d ^{j} \left( \dfrac{e ^{ \gamma ^u_i} e_i^j}{\sum^{}_{t_k \in \pmb{t}} j e ^{\gamma ^u _k}w^j_k} \right)\]

where

\[\gamma^u_i = \Psi \left( \sum^{M}_{k=1} \alpha^0_k + \alpha^u+_k \right)\]

where \(\Psi\) is the digamma function and the expected value of
nucleotide fractions can be expressed as
\[
$$
\begin{align*}
	\mathbb{E} \left( \eta_i \right) &= \dfrac{\alpha_i^0 + \alpha_i ^{\prime} }{\sum^{}_{j} \alpha^0_j + a ^{\prime} _j}= \dfrac{a_i^0 + \alpha_i ^{\prime} }{\hat \alpha^0 +N}\\
	\hat \alpha^0 &= \sum^{M}_{i = 1} \alpha_i^0
\end{align*}
$$
\]
\subsection{Datasets}
\section{Methods}
\subsection{Combining the traditional EM algorithm with the VBEM algorithm}
Since we are only interested in optimizing abundance estimation in the offline phase, we opted to directly change Salmon v1.10.3’s 
source code. We introduce a new flag, $\mathtt{--useBoth}$, which forces the program to run the offline phase twice: 
once with the traditional EM algorithm and once with the VBEM algorithm. 
The program then computes the average of the two different estimates for the $\pmb{\alpha}$ vector before computing its final 
prediction for $\pmb{\eta}$.

\subsection{Running the new implementation on the polyester dataset}
The commands used for the quantification of each sample using EM, VBEM, or both inference strategies, are provided in Fig \ref{code}.
\begin{figure*}[!t]%
    \centering
    {        
        \begin{verbatim}
        salmon quant -i human_index -l IU -1 sample_1.fa.gz -2 sample_2.fa.gz -p 6 --gcBias --useEM \
            --validateMappings -o sample
        salmon quant -i human_index -l IU -1 sample_1.fa.gz -2 sample_2.fa.gz -p 6 --gcBias --vbPrior \
            1e{from -5 to 1} --validateMappings -o sample
        salmon quant -i human_index -l IU -1 sample_1.fa.gz -2 sample_2.fa.gz -p 6 --gcBias --useBoth \
            --vbPrior 1e{from -5 to 1} --validateMappings -o sample
    \end{verbatim}
    }
    \caption{Commands used for the quantification of each sample using EM, VBEM, or both inference strategies, respectively.}
    \label{code}
\end{figure*}

Quantification was run on the 24 samples through a slurm batch script. 
For each run, 5GB of memory was allocated on the CBCB’s Nexus partition. 
In total, every sample had 15 runs, each with different specifications for the inference method (1 EM, 7 VBEM, and 7 combo) 
and the prior value if applicable ($10^{-5}$ to $10^1$).  All other parameters were kept the same among the 360 total runs: 
\begin{enumerate}
    \item We kept the selective alignment feature on for this experiment. This strategy is currently the default for the pipeline, 
    and it allows Salmon to adopt a more sensitive approach during quasi-mapping.
    \item In recent versions of Salmon, the library type can be automatically detected. Nevertheless, we elected to 
    specify the library type for all the runs.
    \item The $\mathtt{–gcBias}$ flag was passed to the pipeline as well, which allows Salmon to identify and  remedy for 
    GC biases. This extra step is essential for the analysis of our synthetic data since they were simulated from an already 
    existing GC bias profile \cite{love_swimming_2018}.
    
\end{enumerate}

\section{Experiment}

We began by editing Salmon’s source code. We wrote a bash script that downloaded our datasets. We also had a script that ran our samples using EM, VBEM, and a combination of the VBEM and EM algorithms. 
First we ran unaltered Salmon using EM and VBEM. For each pass that used VBEM, we specified a prior 10-5 through 101. 
	We then ran the same datasets using the average of both the EM and VBEM algorithms. We ran the EM and VBEM algorithms on the same data, producing an array. FOr each index in this array, the data was averaged and inserted into a new array that was used the final calculations. We again used priors of 10-5 through 101 when considering the VBEM algorithm. 
	We ran this on [Nathalie’s laptop]
	Finally, we compared the results of these passes using the ground truth of the synthetic data. We calculate the MARD of our results. We also calculated the Spearman Correlation and a two-sided Mann-Whitney U test. 

\section{Results}
We encountered different results depending on the prior selected. Overall, we can separate the outcome of combining both inference algorithms into two categories. The first of such categories includes runs with a prior between $10^{-5}$ and $10^1$, inclusive. As shown in figures \ref{speartable} and \ref{mardtable}, both the $\mathrm{MARD}$ values and the Spearman correlation values are lower for estimates inferred from VBEM than for estimates inferred from both methods (Mann-Whitney U test, $P\leq 1.436 \cdot  10^{-7}$ for MARD, $P < 6.202\cdot 10^{-14}$ for Spearman). However, MARD values tended to be similar between estimates inferred from EM and estimates inferred from both methods (Mann-Whitney U test, $P\geq 0.4803$), whereas the Spearman correlation values were slightly higher for EM (Mann-Whitney U test, $P\leq 0.03941$). Because the MARD metrics are negatively impacted by false discovery rates \cite{patro_salmon_2017}, we hypothesize that the results gathered from running both inference algorithms is affected by EM’s tendency to collect more non-zero abundances than VBEM at low prior values. 

This can also be concluded from the analysis of our second category, which includes runs with a prior of 1 or 10. Unlike with the first category, the $\mathrm{MARD}$ and the Spearman correlation values are strikingly lower for the estimates inferred from EM (Mann-Whitney U test, $P < 6.202\cdot 10^{-14}$ for $\mathrm{MARD}$, $P < 6.202\cdot 10^{-14}$ for Spearman), whereas the difference in MARD between VBEM and the combination of both methods was insignificant for runs with a prior of 1 (Mann-Whitney U test, $P = 0.8944$). The difference between the Spearman correlation values and between the MARD values for runs with a prior of 10 was still significant (Mann-Whitney U test, $P = 5.675 \cdot 10^{-11}$ for $\mathrm{MARD}$, $P < 6.202\cdot 10^{-14}$ for Spearman). We predict that this would suggest that for high prior values, the negative impact from false discovery rates by the VBEM algorithm is offsetted by the more accurate abundance estimates from the EM algorithm. 
\begin{table*}[t]
    \caption{Spearman correlation table for first four samples and their subsamples.
    \label{speartable}}
    \tabcolsep=0pt%%
    \begin{tabular*}{\textwidth}{@{\extracolsep{\fill}}lcccccc@{\extracolsep{\fill}}}
    \toprule
    \begin{tabular}{lrrrrrrrr}
        \textbf{} &
          \multicolumn{1}{l}{Sample $1_{01}$} &
          \multicolumn{1}{l}{Sample $1_{02}$} &
          \multicolumn{1}{l}{Sample $2_{01}$} &
          \multicolumn{1}{l}{Sample $2_{02}$} &
          \multicolumn{1}{l}{Sample $3_{01}$} &
          \multicolumn{1}{l}{Sample $3_{02}$} &
          \multicolumn{1}{l}{Sample $4_{01}$} &
          \multicolumn{1}{l}{Sample $4_{02}$} \\
        EM           & 0.932 & 0.932 & 0.933 & 0.932 & 0.937 & 0.936 & 0.934 & 0.933 \\
        both - $10^{-5}$ & 0.933 & 0.933 & 0.934 & 0.933 & 0.939 & 0.938 & 0.935 & 0.934 \\
        VBEM - $10^{-5}$ & 0.952 & 0.951 & 0.953 & 0.951 & 0.957 & 0.956 & 0.954 & 0.952 \\
        both - $10^{-4}$ & 0.933 & 0.933 & 0.934 & 0.933 & 0.939 & 0.938 & 0.935 & 0.934 \\
        VBEM - $10^{-4}$ & 0.952 & 0.951 & 0.953 & 0.951 & 0.957 & 0.956 & 0.954 & 0.952 \\
        both - $10^{-3}$ & 0.933 & 0.933 & 0.934 & 0.933 & 0.939 & 0.938 & 0.935 & 0.934 \\
        VBEM - $10^{-3}$ & 0.952 & 0.951 & 0.953 & 0.951 & 0.957 & 0.956 & 0.954 & 0.952 \\
        both - $10^{-2}$ & 0.933 & 0.933 & 0.934 & 0.933 & 0.939 & 0.938 & 0.935 & 0.934 \\
        VBEM - $10^{-2}$ & 0.952 & 0.951 & 0.953 & 0.952 & 0.957 & 0.956 & 0.954 & 0.952 \\
        both - $10^{-1}$ & 0.933 & 0.933 & 0.934 & 0.933 & 0.939 & 0.937 & 0.935 & 0.934 \\
        VBEM - $10^{-1}$ & 0.951 & 0.950 & 0.953 & 0.951 & 0.957 & 0.955 & 0.954 & 0.951 \\
        both - $10^{0}$  & 0.755 & 0.754 & 0.755 & 0.754 & 0.755 & 0.755 & 0.755 & 0.754 \\
        VBEM - $10^{0}$  & 0.751 & 0.751 & 0.751 & 0.751 & 0.752 & 0.752 & 0.751 & 0.751 \\
        both - $10^{1}$  & 0.736 & 0.735 & 0.736 & 0.736 & 0.738 & 0.737 & 0.736 & 0.736 \\
        VBEM - $10^{1}$  & 0.715 & 0.715 & 0.715 & 0.716 & 0.717 & 0.717 & 0.715 & 0.715 \\
        \botrule    
    \end{tabular}
    
    \end{tabular*}
    \end{table*}
\section{Statistical Evaluation}
The statistical evaluation of the different inference algorithms was performed with R. 
Because only the original simulated read counts were made available to, we conducted our comparisons on the estimated number of reads 
given by Salmon. We chose two statistical metrics for our analysis:
The mean absolute relative difference, or $\mathrm{MARD}$, 
was already introduced in the original Salmon publication. 
It is the average of the absolute relative difference $\mathrm{ARD}$ for each transcript $i$, 
where $x_i$ is the original simulated read count, and $y_i$ is the read count reported by Salmon.
\[
    \mathrm{ARD}_i = \begin{cases}
        0 &\quad\text{if}\quad x_i = y_i = 0\\
        \dfrac{\left| x_i-y_i \right| }{x_i-y_i} &\quad\text{otherwise}\quad
    \end{cases}
\]
\[
    \mathrm{MARD} = \frac{1}{M} \sum^{M}_{i = 1}\mathrm{ARD}_i
\]
The Spearman correlation, also described in the original publication, calculates the relationship between the rank for our results and the rank for the simulated data. It was computed through the core.test function where we indicated a two sided alternative hypothesis.


\subsection{Spearman Correlation}
\[
    \rho = 1 - \dfrac{6 \sum{d^2_i}}{n(n^2-1)}
\]
\subsection{Mann–Whitney U Test}

MATH
Additionally, we were also required to measure the differences between the three inference algorithms. 
We used the Mann-Whitney test (also known as the Wilcoxon test) 
to calculate the p value for such analysis, and computations were also performed in R 
with the wilcox.test function

\section{Discussion}

We were unable to expand on our project with the time given but in future work, we would like to expand upon ways to best combine the results of the EM and VBEM algorithms. Currently we are only taking the mean of the two results. However, we would also like to experiment with the weighted average of these two algorithms also taking into account various prior sizes. We would like to see if this also improves the estimates of Salmon. 

We would also like to verify the time it takes to run our combined approach. We would like to verify that running this approach will not greatly impact the runtime of Salmon. While we expect the runtime to increase; however, we would like to verify that it is not detrimental to the use of Salmon. 

In future work, it will be beneficial to get a measure of uncertainty when running our trials. We were unable to run this during our initial experiments. However, we would like to know if there is a level of uncertainty with our results when we use the average of the EM and VBEM algorithms. We would like to know if this uncertainty increases or decreases, and if it does so significantly. 



% \section{Introduction}\label{sec1}

% The Introduction section, of referenced text \cite{bib1} expands on the background of the work (some overlap with the Abstract is acceptable). The introduction should not include subheadings.

% Springer Nature does not impose a strict layout as standard however authors are advised to check the individual requirements for the journal they are planning to submit to as there may be journal-level preferences. When preparing your text please also be aware that some stylistic choices are not supported in full text XML (publication version), including coloured font. These will not be replicated in the typeset article if it is accepted. 

% \section{Results}\label{sec2}

% Sample body text. Sample body text. Sample body text. Sample body text. Sample body text. Sample body text. Sample body text. Sample body text.

% \section{This is an example for first level head---section head}\label{sec3}

% \subsection{This is an example for second level head---subsection head}\label{subsec2}

% \subsubsection{This is an example for third level head---subsubsection head}\label{subsubsec2}

% Sample body text. Sample body text. Sample body text. Sample body text. Sample body text. Sample body text. Sample body text. Sample body text. 

% \section{Equations}\label{sec4}

% Equations in \LaTeX\ can either be inline or on-a-line by itself (``display equations''). For
% inline equations use the \verb+$...$+ commands. E.g.: The equation
% $H\psi = E \psi$ is written via the command \verb+$H \psi = E \psi$+.

% For display equations (with auto generated equation numbers)
% one can use the equation or align environments:
% \begin{equation}
% \|\tilde{X}(k)\|^2 \leq\frac{\sum\limits_{i=1}^{p}\left\|\tilde{Y}_i(k)\right\|^2+\sum\limits_{j=1}^{q}\left\|\tilde{Z}_j(k)\right\|^2 }{p+q}.\label{eq1}
% \end{equation}
% where,
% \begin{align}
% D_\mu &=  \partial_\mu - ig \frac{\lambda^a}{2} A^a_\mu \nonumber \\
% F^a_{\mu\nu} &= \partial_\mu A^a_\nu - \partial_\nu A^a_\mu + g f^{abc} A^b_\mu A^a_\nu \label{eq2}
% \end{align}
% Notice the use of \verb+\nonumber+ in the align environment at the end
% of each line, except the last, so as not to produce equation numbers on
% lines where no equation numbers are required. The \verb+\label{}+ command
% should only be used at the last line of an align environment where
% \verb+\nonumber+ is not used.
% \begin{equation}
% Y_\infty = \left( \frac{m}{\textrm{GeV}} \right)^{-3}
%     \left[ 1 + \frac{3 \ln(m/\textrm{GeV})}{15}
%     + \frac{\ln(c_2/5)}{15} \right]
% \end{equation}
% The class file also supports the use of \verb+\mathbb{}+, \verb+\mathscr{}+ and
% \verb+\mathcal{}+ commands. As such \verb+\mathbb{R}+, \verb+\mathscr{R}+
% and \verb+\mathcal{R}+ produces $\mathbb{R}$, $\mathscr{R}$ and $\mathcal{R}$
% respectively (refer Subsubsection~\ref{subsubsec2}).

% \section{Tables}\label{sec5}

% Tables can be inserted via the normal table and tabular environment. To put
% footnotes inside tables you should use \verb+\footnotetext[]{...}+ tag.
% The footnote appears just below the table itself (refer Tables~\ref{tab1} and \ref{tab2}). 
% For the corresponding footnotemark use \verb+\footnotemark[...]+

% \begin{table}[h]
% \caption{Caption text}\label{tab1}%
% \begin{tabular}{@{}llll@{}}
% \toprule
% Column 1 & Column 2  & Column 3 & Column 4\\
% \midrule
% row 1    & data 1   & data 2  & data 3  \\
% row 2    & data 4   & data 5\footnotemark[1]  & data 6  \\
% row 3    & data 7   & data 8  & data 9\footnotemark[2]  \\
% \botrule
% \end{tabular}
% \footnotetext{Source: This is an example of table footnote. This is an example of table footnote.}
% \footnotetext[1]{Example for a first table footnote. This is an example of table footnote.}
% \footnotetext[2]{Example for a second table footnote. This is an example of table footnote.}
% \end{table}

% \noindent
% The input format for the above table is as follows:

% %%=============================================%%
% %% For presentation purpose, we have included  %%
% %% \bigskip command. Please ignore this.       %%
% %%=============================================%%
% \bigskip
% \begin{verbatim}
% \begin{table}[<placement-specifier>]
% \caption{<table-caption>}\label{<table-label>}%
% \begin{tabular}{@{}llll@{}}
% \toprule
% Column 1 & Column 2 & Column 3 & Column 4\\
% \midrule
% row 1 & data 1 & data 2	 & data 3 \\
% row 2 & data 4 & data 5\footnotemark[1] & data 6 \\
% row 3 & data 7 & data 8	 & data 9\footnotemark[2]\\
% \botrule
% \end{tabular}
% \footnotetext{Source: This is an example of table footnote. 
% This is an example of table footnote.}
% \footnotetext[1]{Example for a first table footnote.
% This is an example of table footnote.}
% \footnotetext[2]{Example for a second table footnote. 
% This is an example of table footnote.}
% \end{table}
% \end{verbatim}
% \bigskip
% %%=============================================%%
% %% For presentation purpose, we have included  %%
% %% \bigskip command. Please ignore this.       %%
% %%=============================================%%

% \begin{table}[h]
% \caption{Example of a lengthy table which is set to full textwidth}\label{tab2}
% \begin{tabular*}{\textwidth}{@{\extracolsep\fill}lcccccc}
% \toprule%
% & \multicolumn{3}{@{}c@{}}{Element 1\footnotemark[1]} & \multicolumn{3}{@{}c@{}}{Element 2\footnotemark[2]} \\\cmidrule{2-4}\cmidrule{5-7}%
% Project & Energy & $\sigma_{calc}$ & $\sigma_{expt}$ & Energy & $\sigma_{calc}$ & $\sigma_{expt}$ \\
% \midrule
% Element 3  & 990 A & 1168 & $1547\pm12$ & 780 A & 1166 & $1239\pm100$\\
% Element 4  & 500 A & 961  & $922\pm10$  & 900 A & 1268 & $1092\pm40$\\
% \botrule
% \end{tabular*}
% \footnotetext{Note: This is an example of table footnote. This is an example of table footnote this is an example of table footnote this is an example of~table footnote this is an example of table footnote.}
% \footnotetext[1]{Example for a first table footnote.}
% \footnotetext[2]{Example for a second table footnote.}
% \end{table}

% In case of double column layout, tables which do not fit in single column width should be set to full text width. For this, you need to use \verb+\begin{table*}+ \verb+...+ \verb+\end{table*}+ instead of \verb+\begin{table}+ \verb+...+ \verb+\end{table}+ environment. Lengthy tables which do not fit in textwidth should be set as rotated table. For this, you need to use \verb+\begin{sidewaystable}+ \verb+...+ \verb+\end{sidewaystable}+ instead of \verb+\begin{table*}+ \verb+...+ \verb+\end{table*}+ environment. This environment puts tables rotated to single column width. For tables rotated to double column width, use \verb+\begin{sidewaystable*}+ \verb+...+ \verb+\end{sidewaystable*}+.

% \begin{sidewaystable}
% \caption{Tables which are too long to fit, should be written using the ``sidewaystable'' environment as shown here}\label{tab3}
% \begin{tabular*}{\textheight}{@{\extracolsep\fill}lcccccc}
% \toprule%
% & \multicolumn{3}{@{}c@{}}{Element 1\footnotemark[1]}& \multicolumn{3}{@{}c@{}}{Element\footnotemark[2]} \\\cmidrule{2-4}\cmidrule{5-7}%
% Projectile & Energy	& $\sigma_{calc}$ & $\sigma_{expt}$ & Energy & $\sigma_{calc}$ & $\sigma_{expt}$ \\
% \midrule
% Element 3 & 990 A & 1168 & $1547\pm12$ & 780 A & 1166 & $1239\pm100$ \\
% Element 4 & 500 A & 961  & $922\pm10$  & 900 A & 1268 & $1092\pm40$ \\
% Element 5 & 990 A & 1168 & $1547\pm12$ & 780 A & 1166 & $1239\pm100$ \\
% Element 6 & 500 A & 961  & $922\pm10$  & 900 A & 1268 & $1092\pm40$ \\
% \botrule
% \end{tabular*}
% \footnotetext{Note: This is an example of table footnote this is an example of table footnote this is an example of table footnote this is an example of~table footnote this is an example of table footnote.}
% \footnotetext[1]{This is an example of table footnote.}
% \end{sidewaystable}

% \section{Figures}\label{sec6}

% As per the \LaTeX\ standards you need to use eps images for \LaTeX\ compilation and \verb+pdf/jpg/png+ images for \verb+PDFLaTeX+ compilation. This is one of the major difference between \LaTeX\ and \verb+PDFLaTeX+. Each image should be from a single input .eps/vector image file. Avoid using subfigures. The command for inserting images for \LaTeX\ and \verb+PDFLaTeX+ can be generalized. The package used to insert images in \verb+LaTeX/PDFLaTeX+ is the graphicx package. Figures can be inserted via the normal figure environment as shown in the below example:

% %%=============================================%%
% %% For presentation purpose, we have included  %%
% %% \bigskip command. Please ignore this.       %%
% %%=============================================%%
% \bigskip
% \begin{verbatim}
% \begin{figure}[<placement-specifier>]
% \centering
% \includegraphics{<eps-file>}
% \caption{<figure-caption>}\label{<figure-label>}
% \end{figure}
% \end{verbatim}
% \bigskip
% %%=============================================%%
% %% For presentation purpose, we have included  %%
% %% \bigskip command. Please ignore this.       %%
% %%=============================================%%

% \begin{figure}[h]
% \centering
% \includegraphics[width=0.9\textwidth]{fig.eps}
% \caption{This is a widefig. This is an example of long caption this is an example of long caption  this is an example of long caption this is an example of long caption}\label{fig1}
% \end{figure}

% In case of double column layout, the above format puts figure captions/images to single column width. To get spanned images, we need to provide \verb+\begin{figure*}+ \verb+...+ \verb+\end{figure*}+.

% For sample purpose, we have included the width of images in the optional argument of \verb+\includegraphics+ tag. Please ignore this. 

% \section{Algorithms, Program codes and Listings}\label{sec7}

% Packages \verb+algorithm+, \verb+algorithmicx+ and \verb+algpseudocode+ are used for setting algorithms in \LaTeX\ using the format:

% %%=============================================%%
% %% For presentation purpose, we have included  %%
% %% \bigskip command. Please ignore this.       %%
% %%=============================================%%
% \bigskip
% \begin{verbatim}
% \begin{algorithm}
% \caption{<alg-caption>}\label{<alg-label>}
% \begin{algorithmic}[1]
% . . .
% \end{algorithmic}
% \end{algorithm}
% \end{verbatim}
% \bigskip
% %%=============================================%%
% %% For presentation purpose, we have included  %%
% %% \bigskip command. Please ignore this.       %%
% %%=============================================%%

% You may refer above listed package documentations for more details before setting \verb+algorithm+ environment. For program codes, the ``verbatim'' package is required and the command to be used is \verb+\begin{verbatim}+ \verb+...+ \verb+\end{verbatim}+. 

% Similarly, for \verb+listings+, use the \verb+listings+ package. \verb+\begin{lstlisting}+ \verb+...+ \verb+\end{lstlisting}+ is used to set environments similar to \verb+verbatim+ environment. Refer to the \verb+lstlisting+ package documentation for more details.

% A fast exponentiation procedure:

% \lstset{texcl=true,basicstyle=\small\sf,commentstyle=\small\rm,mathescape=true,escapeinside={(*}{*)}}
% \begin{lstlisting}
% begin
%   for $i:=1$ to $10$ step $1$ do
%       expt($2,i$);  
%       newline() od                (*\textrm{Comments will be set flush to the right margin}*)
% where
% proc expt($x,n$) $\equiv$
%   $z:=1$;
%   do if $n=0$ then exit fi;
%      do if odd($n$) then exit fi;                 
%         comment: (*\textrm{This is a comment statement;}*)
%         $n:=n/2$; $x:=x*x$ od;
%      { $n>0$ };
%      $n:=n-1$; $z:=z*x$ od;
%   print($z$). 
% end
% \end{lstlisting}

% \begin{algorithm}
% \caption{Calculate $y = x^n$}\label{algo1}
% \begin{algorithmic}[1]
% \Require $n \geq 0 \vee x \neq 0$
% \Ensure $y = x^n$ 
% \State $y \Leftarrow 1$
% \If{$n < 0$}\label{algln2}
%         \State $X \Leftarrow 1 / x$
%         \State $N \Leftarrow -n$
% \Else
%         \State $X \Leftarrow x$
%         \State $N \Leftarrow n$
% \EndIf
% \While{$N \neq 0$}
%         \If{$N$ is even}
%             \State $X \Leftarrow X \times X$
%             \State $N \Leftarrow N / 2$
%         \Else[$N$ is odd]
%             \State $y \Leftarrow y \times X$
%             \State $N \Leftarrow N - 1$
%         \EndIf
% \EndWhile
% \end{algorithmic}
% \end{algorithm}

% %%=============================================%%
% %% For presentation purpose, we have included  %%
% %% \bigskip command. Please ignore this.       %%
% %%=============================================%%
% \bigskip
% \begin{minipage}{\hsize}%
% \lstset{frame=single,framexleftmargin=-1pt,framexrightmargin=-17pt,framesep=12pt,linewidth=0.98\textwidth,language=pascal}% Set your language (you can change the language for each code-block optionally)
% %%% Start your code-block
% \begin{lstlisting}
% for i:=maxint to 0 do
% begin
% { do nothing }
% end;
% Write('Case insensitive ');
% Write('Pascal keywords.');
% \end{lstlisting}
% \end{minipage}

% \section{Cross referencing}\label{sec8}

% Environments such as figure, table, equation and align can have a label
% declared via the \verb+\label{#label}+ command. For figures and table
% environments use the \verb+\label{}+ command inside or just
% below the \verb+\caption{}+ command. You can then use the
% \verb+\ref{#label}+ command to cross-reference them. As an example, consider
% the label declared for Figure~\ref{fig1} which is
% \verb+\label{fig1}+. To cross-reference it, use the command 
% \verb+Figure \ref{fig1}+, for which it comes up as
% ``Figure~\ref{fig1}''. 

% To reference line numbers in an algorithm, consider the label declared for the line number 2 of Algorithm~\ref{algo1} is \verb+\label{algln2}+. To cross-reference it, use the command \verb+\ref{algln2}+ for which it comes up as line~\ref{algln2} of Algorithm~\ref{algo1}.

% \subsection{Details on reference citations}\label{subsec7}

% Standard \LaTeX\ permits only numerical citations. To support both numerical and author-year citations this template uses \verb+natbib+ \LaTeX\ package. For style guidance please refer to the template user manual.

% Here is an example for \verb+\cite{...}+: \cite{bib1}. Another example for \verb+\citep{...}+: \citep{bib2}. For author-year citation mode, \verb+\cite{...}+ prints Jones et al. (1990) and \verb+\citep{...}+ prints (Jones et al., 1990).

% All cited bib entries are printed at the end of this article: \cite{bib3}, \cite{bib4}, \cite{bib5}, \cite{bib6}, \cite{bib7}, \cite{bib8}, \cite{bib9}, \cite{bib10}, \cite{bib11}, \cite{bib12} and \cite{bib13}.


% \section{Examples for theorem like environments}\label{sec10}

% For theorem like environments, we require \verb+amsthm+ package. There are three types of predefined theorem styles exists---\verb+thmstyleone+, \verb+thmstyletwo+ and \verb+thmstylethree+ 

% %%=============================================%%
% %% For presentation purpose, we have included  %%
% %% \bigskip command. Please ignore this.       %%
% %%=============================================%%
% \bigskip
% \begin{tabular}{|l|p{19pc}|}
% \hline
% \verb+thmstyleone+ & Numbered, theorem head in bold font and theorem text in italic style \\\hline
% \verb+thmstyletwo+ & Numbered, theorem head in roman font and theorem text in italic style \\\hline
% \verb+thmstylethree+ & Numbered, theorem head in bold font and theorem text in roman style \\\hline
% \end{tabular}
% \bigskip
% %%=============================================%%
% %% For presentation purpose, we have included  %%
% %% \bigskip command. Please ignore this.       %%
% %%=============================================%%

% For mathematics journals, theorem styles can be included as shown in the following examples:

% \begin{theorem}[Theorem subhead]\label{thm1}
% Example theorem text. Example theorem text. Example theorem text. Example theorem text. Example theorem text. 
% Example theorem text. Example theorem text. Example theorem text. Example theorem text. Example theorem text. 
% Example theorem text. 
% \end{theorem}

% Sample body text. Sample body text. Sample body text. Sample body text. Sample body text. Sample body text. Sample body text. Sample body text.

% \begin{proposition}
% Example proposition text. Example proposition text. Example proposition text. Example proposition text. Example proposition text. 
% Example proposition text. Example proposition text. Example proposition text. Example proposition text. Example proposition text. 
% \end{proposition}

% Sample body text. Sample body text. Sample body text. Sample body text. Sample body text. Sample body text. Sample body text. Sample body text.

% \begin{example}
% Phasellus adipiscing semper elit. Proin fermentum massa
% ac quam. Sed diam turpis, molestie vitae, placerat a, molestie nec, leo. Maecenas lacinia. Nam ipsum ligula, eleifend
% at, accumsan nec, suscipit a, ipsum. Morbi blandit ligula feugiat magna. Nunc eleifend consequat lorem. 
% \end{example}

% Sample body text. Sample body text. Sample body text. Sample body text. Sample body text. Sample body text. Sample body text. Sample body text.

% \begin{remark}
% Phasellus adipiscing semper elit. Proin fermentum massa
% ac quam. Sed diam turpis, molestie vitae, placerat a, molestie nec, leo. Maecenas lacinia. Nam ipsum ligula, eleifend
% at, accumsan nec, suscipit a, ipsum. Morbi blandit ligula feugiat magna. Nunc eleifend consequat lorem. 
% \end{remark}

% Sample body text. Sample body text. Sample body text. Sample body text. Sample body text. Sample body text. Sample body text. Sample body text.

% \begin{definition}[Definition sub head]
% Example definition text. Example definition text. Example definition text. Example definition text. Example definition text. Example definition text. Example definition text. Example definition text. 
% \end{definition}

% Additionally a predefined ``proof'' environment is available: \verb+\begin{proof}+ \verb+...+ \verb+\end{proof}+. This prints a ``Proof'' head in italic font style and the ``body text'' in roman font style with an open square at the end of each proof environment. 

% \begin{proof}
% Example for proof text. Example for proof text. Example for proof text. Example for proof text. Example for proof text. Example for proof text. Example for proof text. Example for proof text. Example for proof text. Example for proof text. 
% \end{proof}

% Sample body text. Sample body text. Sample body text. Sample body text. Sample body text. Sample body text. Sample body text. Sample body text.

% \begin{proof}[Proof of Theorem~{\upshape\ref{thm1}}]
% Example for proof text. Example for proof text. Example for proof text. Example for proof text. Example for proof text. Example for proof text. Example for proof text. Example for proof text. Example for proof text. Example for proof text. 
% \end{proof}

% \noindent
% For a quote environment, use \verb+\begin{quote}...\end{quote}+
% \begin{quote}
% Quoted text example. Aliquam porttitor quam a lacus. Praesent vel arcu ut tortor cursus volutpat. In vitae pede quis diam bibendum placerat. Fusce elementum
% convallis neque. Sed dolor orci, scelerisque ac, dapibus nec, ultricies ut, mi. Duis nec dui quis leo sagittis commodo.
% \end{quote}

% Sample body text. Sample body text. Sample body text. Sample body text. Sample body text (refer Figure~\ref{fig1}). Sample body text. Sample body text. Sample body text (refer Table~\ref{tab3}). 

% \section{Methods}\label{sec11}

% Topical subheadings are allowed. Authors must ensure that their Methods section includes adequate experimental and characterization data necessary for others in the field to reproduce their work. Authors are encouraged to include RIIDs where appropriate. 

% \textbf{Ethical approval declarations} (only required where applicable) Any article reporting experiment/s carried out on (i)~live vertebrate (or higher invertebrates), (ii)~humans or (iii)~human samples must include an unambiguous statement within the methods section that meets the following requirements: 

% \begin{enumerate}[1.]
% \item Approval: a statement which confirms that all experimental protocols were approved by a named institutional and/or licensing committee. Please identify the approving body in the methods section

% \item Accordance: a statement explicitly saying that the methods were carried out in accordance with the relevant guidelines and regulations

% \item Informed consent (for experiments involving humans or human tissue samples): include a statement confirming that informed consent was obtained from all participants and/or their legal guardian/s
% \end{enumerate}

% If your manuscript includes potentially identifying patient/participant information, or if it describes human transplantation research, or if it reports results of a clinical trial then  additional information will be required. Please visit (\url{https://www.nature.com/nature-research/editorial-policies}) for Nature Portfolio journals, (\url{https://www.springer.com/gp/authors-editors/journal-author/journal-author-helpdesk/publishing-ethics/14214}) for Springer Nature journals, or (\url{https://www.biomedcentral.com/getpublished/editorial-policies\#ethics+and+consent}) for BMC.

% \section{Discussion}\label{sec12}

% Discussions should be brief and focused. In some disciplines use of Discussion or `Conclusion' is interchangeable. It is not mandatory to use both. Some journals prefer a section `Results and Discussion' followed by a section `Conclusion'. Please refer to Journal-level guidance for any specific requirements. 

% \section{Conclusion}\label{sec13}

% Conclusions may be used to restate your hypothesis or research question, restate your major findings, explain the relevance and the added value of your work, highlight any limitations of your study, describe future directions for research and recommendations. 

% In some disciplines use of Discussion or 'Conclusion' is interchangeable. It is not mandatory to use both. Please refer to Journal-level guidance for any specific requirements. 

% \backmatter

% \bmhead{Supplementary information}

% If your article has accompanying supplementary file/s please state so here. 

% Authors reporting data from electrophoretic gels and blots should supply the full unprocessed scans for key as part of their Supplementary information. This may be requested by the editorial team/s if it is missing.

% Please refer to Journal-level guidance for any specific requirements.

% \bmhead{Acknowledgements}

% Acknowledgements are not compulsory. Where included they should be brief. Grant or contribution numbers may be acknowledged.

% Please refer to Journal-level guidance for any specific requirements.

% \section*{Declarations}

% Some journals require declarations to be submitted in a standardised format. Please check the Instructions for Authors of the journal to which you are submitting to see if you need to complete this section. If yes, your manuscript must contain the following sections under the heading `Declarations':

% \begin{itemize}
% \item Funding
% \item Conflict of interest/Competing interests (check journal-specific guidelines for which heading to use)
% \item Ethics approval and consent to participate
% \item Consent for publication
% \item Data availability 
% \item Materials availability
% \item Code availability 
% \item Author contribution
% \end{itemize}

% \noindent
% If any of the sections are not relevant to your manuscript, please include the heading and write `Not applicable' for that section. 

% %%===================================================%%
% %% For presentation purpose, we have included        %%
% %% \bigskip command. Please ignore this.             %%
% %%===================================================%%
% \bigskip
% \begin{flushleft}%
% Editorial Policies for:

% \bigskip\noindent
% Springer journals and proceedings: \url{https://www.springer.com/gp/editorial-policies}

% \bigskip\noindent
% Nature Portfolio journals: \url{https://www.nature.com/nature-research/editorial-policies}

% \bigskip\noindent
% \textit{Scientific Reports}: \url{https://www.nature.com/srep/journal-policies/editorial-policies}

% \bigskip\noindent
% BMC journals: \url{https://www.biomedcentral.com/getpublished/editorial-policies}
% \end{flushleft}

% \begin{appendices}

% \section{Section title of first appendix}\label{secA1}

% An appendix contains supplementary information that is not an essential part of the text itself but which may be helpful in providing a more comprehensive understanding of the research problem or it is information that is too cumbersome to be included in the body of the paper.

% %%=============================================%%
% %% For submissions to Nature Portfolio Journals %%
% %% please use the heading ``Extended Data''.   %%
% %%=============================================%%

% %%=============================================================%%
% %% Sample for another appendix section			       %%
% %%=============================================================%%

% %% \section{Example of another appendix section}\label{secA2}%
% %% Appendices may be used for helpful, supporting or essential material that would otherwise 
% %% clutter, break up or be distracting to the text. Appendices can consist of sections, figures, 
% %% tables and equations etc.

% \end{appendices}

%%===========================================================================================%%
%% If you are submitting to one of the Nature Portfolio journals, using the eJP submission   %%
%% system, please include the references within the manuscript file itself. You may do this  %%
%% by copying the reference list from your .bbl file, paste it into the main manuscript .tex %%
%% file, and delete the associated \verb+\bibliography+ commands.                            %%
%%===========================================================================================%%

\bibliography{export-data}% common bib file
%% if required, the content of .bbl file can be included here once bbl is generated
%%\input sn-article.bbl


\end{document}
